\paragraph{The solver.} The free-boundary problem is solved in the log ratio $y$ on $[b,\bar y]$ with $\bar y=8$, by collocation with the threshold $b$ as a free parameter (SciPy's \texttt{solve\_bvp}, tolerance $10^{-7}$, retried at $10^{-6}$ with a finer mesh and more starting values when needed). The boundary conditions are value matching and smooth fit at $b$, $F(b)=F'(b)=0$, and a far-field condition at $\bar y$ taken from the growth of the solution: for large $y$, $F(y)=Ae^y-y/\rho+C+o(1)$ with $A=\alpha/(\rho-D)$, so the solver imposes $F'(\bar y)=Ae^{\bar y}-1/\rho$. The Hamiltonian is coded as $\log(1+\rho\max\{F',0\})$, the form used in Lemma~\OAref{lem:reg}. The solver is a separate implementation from the one in earlier versions of this paper; at the three calibrations of Table~\ref{tab:calib} the two agree to the digits reported. An independent shooting solver, used in the interactive companion, reproduces the same thresholds and welfare losses to four digits.

\paragraph{Checks of Theorem~\ref{thm:main}.} The grid $\rho\in\{.05,.1,.2\}$, $\alpha\in\{.2,1,5\}$, $\lambda\in\{1,2,10\}$, $\delta\in\{-.03,0,.03,.1\}$, $\sigma\in\{.1,.25,.4\}$ has 324 points, of which 243 satisfy $\rho>D$. The problem was solved at all 243. At each I checked $\bar x\le x^*<1$; the lower bound of Lemma~\OAref{lem:above}; and the sign conditions $F\ge0$, $F'\ge0$. All held in every case, and the smallest margin was $x^*-\bar x=0.016$. At $\rho=.1$, $\alpha=1$, $\lambda=2$, $\delta=0$, $\sigma=.25$: $\bar x=0.2633\le x^*=0.57932<1$; the curvature $F''(b+)$ is $26.926$ numerically against $26.924$ from Theorem~\ref{thm:main}(b); and the jump in $F'''$ at $y=0$ is $32.02$ against $32$.

\paragraph{The burst bound.} Inequality \eqref{eq:legendre} of Lemma~\ref{lem:burst} was tested at $2\times10^5$ random triples $(T,s,q)$ with $T\le80$ and $q\in[e^{-12},e^{12}]$, with no violation (largest excess $-1.4\times10^{-11}$). For explicit burst plans the gain term $e^{-\rho T}Q(T)/\rho-\int_0^Te^{-\rho t}H_\rho\,dt$ stayed below $g(T)$ in every case; the most aggressive, a rate of $10^8$ over the last $10^{-4}$ years of a 30-year horizon, reached $1.47$ against $g(30)=1.98$.

\paragraph{Convexity, for Theorem~\ref{thm:vol}(b).} On the same 243 solutions I computed $F''-F'$ on a grid of 20,000 points above the threshold, using \eqref{eq:ode} and, as a check, finite differences. Setting aside three cases in which it dipped below zero only within $0.03$ of the truncation point $\bar y=8$, an artifact of the far-field condition, it was nonnegative throughout for every case with $\lambda=1$, where convexity is proved, and for 143 of the 162 cases with $\lambda>1$. In the other 19 it turned negative in a band around par; all 19 have $\alpha\ge1$ and $\lambda\ge2$. At the three calibrations of Table~\ref{tab:calib} its minimum above the threshold is $5.75$, $14.60$ and $20.00$, the last in the far field, where $F''-F'$ tends to $1/\rho$. For one of them, $(\rho,\alpha,\lambda,\delta,\sigma)=(.05,5,10,-.03,.10)$, I also solved at $\sigma=.15$ with $D$ fixed, by collocation (with $\bar y=8$ and $\bar y=10$) and by shooting. The threshold falls from $0.5416$ to $0.5066$, but $F(0)$ falls from $308.3$ to $288.8$, while $F$ rises near the threshold (at $y=-0.5$, from $24.1$ to $31.2$). So the conclusion $\hat F\ge F$ of Step 1 of the proof of Theorem~\ref{thm:vol} fails there even though the threshold comparison holds.

\paragraph{Result~\ref{res:vol}.} For the grid, $\rho\in\{.05,.1,.2\}$, $\alpha\in\{.2,1,5\}$, $\lambda\in\{1,2,10\}$ and $\sigma\in\{.05,.10,\dots,.40\}$. With the drift gap fixed, $\delta\in\{-.03,0,.03,.1\}$; with the median growth fixed, $\mut\in\{-.05,-.02,0,.02\}$; with the expected growth fixed, $D\in\{-.05,-.02,0,.02\}$. Points with $\rho-D\le0.005$ are excluded. For each parameter combination I solved for $x^*$ at every admissible $\sigma$ and compared consecutive values. Because the proof fails where loss aversion and positional weight are strong, I also searched further out in that direction: with $D$ fixed, $\rho=.05$, $\alpha\in\{1,5,20\}$, $\lambda\in\{10,30,100\}$, $D\in\{-.05,0,.02\}$ and $\sigma\in\{.05,.10,\dots,.30\}$. The threshold fell in all 135 comparisons, by at least $0.005$ each time.

\paragraph{Risky wealth and the raw drift gap.} Proposition~\ref{prop:wealth} transfers Theorem~\ref{thm:vol} to comparisons that hold $D_S$, $\mut_S$ or $\delta_S$ fixed, not the raw drift gap $\delta$. A case shows why the restriction matters: $\rho=.3$, $\alpha=1$, $\lambda=1$, $\delta=0$, $\sigma_R=1$, $\varrho=1$. Raising $\sigma_W$ from $0.8$ to $1.5$ raises $S$ from $0.04$ to $0.25$, but it also lowers $\mut_S$ from $0.18$ to $-0.625$, and the threshold rises from $0.6970$ to $0.8684$ (the same to four digits with $\bar y=6$, $8$ and $10$). The higher volatility of the ratio comes with a much lower median growth, and the second effect dominates.

\paragraph{Welfare.} For an imposed threshold $\tilde x$, the agent's value under optimal consumption solves \eqref{eq:ode} on $(\log\tilde x,\bar y)$ with $-\log(1+\rho F')$ (consumption above $\rho W$ allowed, since an agent held in the race may wish to consume faster to reach the imposed exit), value matching $F(\log\tilde x)=0$, no smooth-fit condition, and the same far-field condition. The agent's optimal share $\rho/(1+\rho F')$ stays bounded in every case computed; its maximum, reached at the imposed exit point, is $1.11\rho$ at the baseline ($\tilde x=0.80$) and $3.04\rho$ in the high-volatility calibration ($\tilde x=0.50$). The loss at par is the difference between the optimal $F(0)$ and this value at $0$. Values and losses are converted to consumption equivalents with log utility, under which a permanent proportional change $\kappa$ in consumption changes $V$ by $\log(1+\kappa)/\rho$. The value $F(0)$ of the race at par is reported as the permanent rise $e^{\rho F(0)}-1$ that would compensate an agent for quitting, and a loss $L$ as the permanent cut $1-e^{-\rho L}$ that would hurt as much.

\paragraph{The cohort illustration.} For Table~\ref{tab:cohort}, each scenario is solved at the baseline preferences, and the cohort's log standing is placed on a grid of 4,001 points spanning five standard deviations either side of its mean, weighted by the normal density. The share out of the race is the weight at or below $b$; the average grind weights $(\rho-q^*(y))/\rho=\rho F'(y)/(1+\rho F'(y))$ above $b$ and zero below. A one-time jump of 10 percent in the benchmark shifts every member's log standing by $-\log 1.1$.

\paragraph{The grind.} The consumption shares in the worked example of Section~\ref{sec:threshold} are $q^*(y)=\rho/(1+\rho F'(y))$ evaluated on the baseline solution; the reported percentages are $(\rho-q^*)/\rho$.

\paragraph{The deadline.} For Table~\ref{tab:deadline}, $F_s$ is computed by backward induction in the time left $s$ on the variational inequality $\max\{-\partial_sF+\tfrac12\sigma^2F''+\mut F'-\rho F-\log(1+\rho\max(F',0))+f,\,-F\}=0$ with $F_0=0$. Each step of $1/100$ year is implicit in the linear part, with the consumption term and the flow taken explicitly, and its complementarity problem ($F\ge0$, with equality where exit is optimal) is solved exactly by the Brennan--Schwartz sweep, the standard method for stopping problems whose stopping region lies below a single boundary. The grid is $y\in[-6,6]$ with step $0.005$, with $F=0$ at the bottom. At the top the condition is $F''-F'=(1-e^{-\rho s})/\rho$: for large $y$ the solution behaves like $A(s)e^y+B(s)y+C(s)$, and the terms in $y$ of the variational inequality give $B'=-\rho B-1$ with $B(0)=0$. As $s\to\infty$ this becomes the far field $Ae^y-y/\rho+C$ of the collocation solver, which is an exact solution of \eqref{eq:ODE} for $y>0$. Moving the top of the grid to $y=8$ changes no reported value, and neither does imposing $F''=F'$ there instead. The boundary $b(s)$ is read off from $\sqrt{F_s}$, which is linear near $b(s)$ because of smooth fit. In all three calibrations the computed $x^*_s$ is nonincreasing in $s$, as part (c) requires; in the two with $\rho>D$ it is within 0.0001 and 0.0008 of the stationary $x^*$ of the collocation solver at $s=100$, as part (d) requires. Halving both steps changes $x^*_s$ by at most 0.002 at every horizon in Table~\ref{tab:deadline}. Code: \texttt{deadline.py}.
